\documentclass{article}
\usepackage{hyperref}
\usepackage{Style}

\nocite{*} % Comentar si quiero citar
%\addbibresource{bibliografia.bib} % Quitar el comentado si quiero usar bibliografia

\begin{document}

\begin{minipage}{2.5cm}
    \includegraphics[width=2cm]{imagen_puc.jpg}
\end{minipage}
\begin{minipage}{12cm}
    {\sc Pontificia Universidad Católica de Chile\\
    Facultad de Matemáticas\\
    Departamento de Matemática\\
    Profesor: Manuel Cabezas -- Ayudante: Benjamín Mateluna}
\end{minipage}
\vspace{2ex}

{\centerline{\bf Introducción a la Geometría - MAT1304}
\centerline{\bf Ayudantía 20}}
\centerline{\bf 02 de junio de 2026}

\vspace{1cm}
\noindent\textbf{Problema 1.} Determine todos los valores posibles de $x_{0}$ tal que la 
intersección entre la recta $2x+y+1=0$ y la recta de ecuación $y=x_{0}x-3$ tenga coordenadas 
enteras.

\vspace{5mm}
\noindent\textbf{Problema 2.} Encuentre la ecuación de la recta que pasa por los puntos de 
intersección de
\begin{equation*}
    \begin{cases}
        x^{2}+y^{2}=4 \\
        (x-1)^{2}+(y-1)^{2}=4
    \end{cases}
\end{equation*}

\vspace{5mm}
\noindent\textbf{Problema 3.} Encuentre el lugar geométrico de todos los puntos que equidistan del
origen $(-1,1)$.

\vspace{5mm}
\begin{dfn}
    Sea $A$ un conjunto. Una función $f:\R\to A$ se dice parametrización de $A$ si es biyectiva.
\end{dfn}

\vspace{2mm}
\noindent\textbf{Problema 4.} Determine una parametrización de los puntos en $\R^{2}$ tales que 
$x^{2}+y^{2}=1$ con $(x,y)\neq(-1,0)$.

\vspace{1cm}
\noindent\textbf{\underline{Ejercicios Propuestos:}}

\vspace{5mm}
\noindent\textbf{Extra 1.} Demuestre que $x^{2}+y^{2}+4x+6y-23=0$ y $x^{2}+y^{2}-8x-10y+25=0$ son
tangentes, esto es, que se intersectan en un único punto.

\vspace{5mm}
\noindent\textbf{Extra 2.} Encuentre todas las soluciones de $a^{2}+b^{2}=c^{2}$ donde 
$a,b,c\in\Z$. \textbf{Hint:} Use el problema 4.

%\printbibliography % Quitar el comentado si quiero usar bibliografia

\end{document}
